\documentclass[12pt]{article}
\usepackage[utf8]{inputenc}
\usepackage[T1]{fontenc}
\usepackage[margin=1in]{geometry}
\usepackage{setspace}
\usepackage{xcolor}
\usepackage[colorlinks=true, urlcolor=blue!60!black, linkcolor=blue!60!black]{hyperref}
\setlength{\parindent}{2em}
\setlength{\parskip}{1.25em}
\renewcommand{\baselinestretch}{1.0}
\begin{document}
\begin{center}
{\LARGE \textit{Virginia v. Black}, 538 U.S. 343 (2003)\par}
\vspace{0.8em}
{\large Case Brief\par}
\end{center}
\vspace{1em}

\section{Memory Jogger}
Reviewing Virginia's cross-burning statute, the Court held that a State may ban cross burning carried out with intent to intimidate, a form of true threat, without violating \textit{R.A.V.}; but four Justices would have held the prima facie provision facially invalid as construed in Barry Black's jury instruction. Seven Justices agreed that Black's conviction could not stand (\href{https://supreme.justia.com/cases/federal/us/538/343/}{Justia}).

\section{Detailed Case Facts}
Virginia Code Ann. \S~18.2-423 (1996) made it a Class 6 felony ``for any person or persons, with the intent of intimidating any person or group of persons, to burn, or cause to be burned, a cross on the property of another, a highway or other public place.'' A second sentence, added in 1968, provided: ``Any such burning of a cross shall be prima facie evidence of an intent to intimidate a person or group of persons.''

Three respondents were convicted in two separate incidents:
\begin{itemize}
  \item \textbf{Barry Black (August 22, 1998):} Black led a Ku Klux Klan rally of 25--30 people on private property in Carroll County, Virginia, with the landowner's permission, visible from a state highway. Speakers made racist and violent remarks, and the rally concluded with the burning of a 25- to 30-foot cross accompanied by ``Amazing Grace.'' At Black's trial, over his First Amendment objection, the court gave Virginia's Model Jury Instruction stating that ``the burning of a cross by itself is sufficient evidence from which you may infer the required intent.'' The jury convicted Black and fined him \$2,500.
  \item \textbf{Richard Elliott and Jonathan O'Mara (May 2, 1998):} Unaffiliated with the Klan, Elliott, O'Mara, and a third man drove onto the Virginia Beach lawn of Elliott's African-American neighbor, James Jubilee, planted a cross, and tried to set it on fire to ``get back'' at Jubilee for complaining about backyard gunfire. O'Mara pleaded guilty to attempted cross burning and conspiracy (reserving his constitutional challenge). Elliott went to trial; his jury received \textit{no} instruction on the prima facie provision and convicted him of attempted cross burning. Both received 90-day sentences (partially suspended for O'Mara) and \$2,500 fines.
\end{itemize}

\section{Procedural History}
After the Virginia Court of Appeals affirmed all three convictions, the Supreme Court of Virginia consolidated the appeals and struck down \S~18.2-423 on its face (262 Va. 764, 553 S.E.2d 738 (2001)). Relying on \textit{R.A.V. v. City of St. Paul}, 505 U.S. 377 (1992), it held that the statute engaged in impermissible content and viewpoint discrimination by singling out cross burning, and ruled in the alternative that the prima facie evidence provision rendered the law unconstitutionally overbroad. The U.S. Supreme Court granted certiorari.

\section{Judicial Votes}
The Supreme Court \textbf{affirmed} as to Barry Black and \textbf{vacated and remanded} as to Richard Elliott and Jonathan O'Mara. The case was argued on December 11, 2002, and decided on April 7, 2003 (\href{https://supreme.justia.com/cases/federal/us/538/343/}{Justia}; \href{https://www.law.cornell.edu/supremecourt/text/538/343}{Cornell LII}). Disposition: 7--2 to invalidate Black's conviction (plurality + Souter/Kennedy/Ginsburg; Scalia agreed the conviction could not stand due to the jury instruction but opposed dismissing the indictment); 5--4 to vacate and remand as to Elliott and O'Mara (plurality + Scalia).

\begin{itemize}
    \item \textbf{Majority Opinion (Parts I--III):} Justice O'Connor delivered the opinion of the Court, joined by Chief Justice Rehnquist and Justices Stevens, Scalia, and Breyer.
    \item \textbf{Plurality Opinion (Parts IV--V):} Justice O'Connor, joined by Chief Justice Rehnquist and Justices Stevens and Breyer.
    \item \textbf{Concurring Opinion:} Justice Stevens filed a concurring opinion.
    \item \textbf{Opinions Concurring in Part and Dissenting in Part:} Justice Scalia filed an opinion concurring in part, concurring in the judgment in part, and dissenting in part, joined by Justice Thomas as to Parts I and II. Justice Souter filed an opinion concurring in the judgment in part and dissenting in part, joined by Justices Kennedy and Ginsburg.
    \item \textbf{Dissenting Opinion:} Justice Thomas filed a dissenting opinion.
\end{itemize}

\section{Holding}
\textbf{Questions Presented:}

\begin{enumerate}
  \item Whether a State may, consistent with the First Amendment and \textit{R.A.V. v. City of St. Paul}, ban cross burning carried out with the intent to intimidate.
  \item Whether Virginia's statutory provision treating any cross burning as ``prima facie evidence of an intent to intimidate'' renders the statute facially unconstitutional under the First Amendment.
\end{enumerate}

\textbf{Holding:}

\begin{enumerate}
  \item \textbf{Core Ban Upheld (Majority, Parts I--III):} Yes. Consistent with the First Amendment and \textit{R.A.V.}, a State may ban cross burning carried out with the intent to intimidate because such intimidation is a ``true threat'' and cross burning is a ``particularly virulent form of intimidation.''
  \item \textbf{Prima Facie Provision Invalidated (Plurality, Parts IV--V):} Yes. As interpreted by Virginia's Model Jury Instruction given in Black's trial, the prima facie evidence provision is facially unconstitutional because it permits conviction based solely on the act of cross burning itself, stripping away contextual factors and chilling core political expression. Black's reversal was affirmed; Elliott's and O'Mara's cases were vacated and remanded for state-court consideration of severability or a narrowing construction.
\end{enumerate}

\section{Analysis of Opinions}
\subsection{Majority Opinion (Parts I--III) and Plurality Opinion (Parts IV--V) (Justice O'Connor)}

\textbf{1. Historical Context of Cross Burning (Part II).} Justice O'Connor traced cross burning's evolution from 14th-century Scottish clan signals to an emblem ``inextricably intertwined with the history of the Ku Klux Klan.'' While the burning cross has functioned both as a tool of racial terror and as a symbol of Klan ideology and group solidarity, when directed at a target to intimidate, ``few if any messages are more powerful.''

\textbf{2. True Threats and Intimidation (Part III-A).} Although the First Amendment protects offensive symbolic conduct (\textit{Texas v. Johnson}, 491 U.S. 397 (1989)), it permits regulation of historic unprotected categories, including ``true threats'' (\textit{Watts v. United States}, 394 U.S. 705 (1969)). True threats encompass serious expressions of intent to commit unlawful violence against a particular individual or group; the speaker need not actually intend to carry out the threat, because banning true threats protects victims from fear of violence and the disruption fear engenders. ``Intimidation in the constitutionally proscribable sense'' is a type of true threat directed with intent to place the victim in fear of bodily harm or death.

\textbf{3. Consistency with \textit{R.A.V. v. St. Paul} (Part III-B).} The Virginia Supreme Court misread \textit{R.A.V.} as barring all content discrimination within unprotected speech. \textit{R.A.V.} expressly recognized an exception when ``the basis for the content discrimination consists entirely of the very reason the entire class of speech at issue is proscribable'' (505 U.S. at 388)---such as banning only the most prurient obscenity or threats against the President. Unlike St. Paul's ordinance, which targeted bias-motivated fighting words only on ``specified disfavored topics'' (race, color, creed, religion, or gender), Virginia's statute banned \textit{all} intimidating cross burning regardless of motive or victim identity. Because cross burning has a uniquely virulent history as a signal of impending violence, Virginia could single it out within the category of true threats.

\textbf{4. Facial Invalidity of the Prima Facie Provision (Parts IV--V, Plurality).} Treating Virginia's unrepudiated Model Jury Instruction as binding state law, the plurality held that the prima facie clause ``strips away the very reason why a State may ban cross burning with the intent to intimidate.'' It permits arrest, prosecution, and conviction whenever a defendant exercises the right not to mount a defense, and skews jury deliberations even when a defense is offered. Because cross burning is not always intimidating---it may be protected political ideology at a rally or artistic expression in films like \textit{Mississippi Burning}---blurring the line without examining context creates an unacceptable risk of suppressing core political speech.

\subsection{Concurring Opinion (Justice Stevens)}

Justice Stevens joined O'Connor's opinion, reiterating his and Justice White's \textit{R.A.V.} concurrences that intimidating cross burning is clearly unprotected regardless of whether the State bans all other threatening conduct.

\subsection{Opinion Concurring in Part, Concurring in the Judgment in Part, and Dissenting in Part (Justice Scalia)}

Justice Scalia joined Parts I--III but attacked the plurality's facial invalidation of the prima facie clause. Under Virginia law, ``prima facie evidence'' merely lets the prosecution survive a motion to strike until rebutted, which fails substantial-overbreadth review because it risks improper conviction only for defendants who burn a cross non-intimidatingly and decline to present a defense. Scalia also faulted the plurality for facially invalidating a statute based on a non-binding model instruction while simultaneously inviting a saving construction on remand. He agreed Black's conviction failed because of the flawed instruction, but would have permitted retrial rather than dismissing the indictment.

\subsection{Opinion Concurring in the Judgment in Part and Dissenting in Part (Justice Souter)}

Justice Souter concluded the entire Virginia statute violated \textit{R.A.V.} Because a burning cross is uniquely tied to white-supremacist ideology, singling it out does not fit \textit{R.A.V.}'s ``virulence'' exception cleanly; moreover, the prima facie clause skews juries toward conviction in equivocal ideological cases, confirming a realistic possibility that ``official suppression of ideas is afoot.'' Souter would have invalidated all three prosecutions without remand.

\subsection{Dissenting Opinion (Justice Thomas)}

Justice Thomas argued that intimidating cross burning is ``conduct, not expression,'' emphasizing that Virginia's segregationist 1952 legislature sought to curb Klan terrorism rather than suppress racist ideas. Even if the First Amendment applied, the prima facie clause operates merely as a constitutionally permissible ``permissive inference'' that leaves the reasonable-doubt burden intact.

\subsection{Subsequent Developments}

\begin{itemize}

\item \textbf{Subjective Mens Rea Required (\textit{Counterman v.\ Colorado}):} In \textit{Counterman} (7--2, opinion by Justice Kagan), the Court held that to provide ``breathing space'' against self-censorship, the First Amendment requires the State in true-threats prosecutions to prove a subjective mens rea of at least \textbf{recklessness}---that the speaker ``consciously disregarded a substantial risk that his communications would be viewed as threatening violence.'' Crucially, the \textit{Counterman} majority clarified (n.~3) that \textit{Black}'s references to ``intent'' reflected Virginia's statutory element rather than a constitutional floor of purpose or knowledge. (\href{https://supreme.justia.com/cases/federal/us/600/22-138/}{\textit{Counterman v.\ Colorado}, 600 U.S.\ 66 (2023)}.)

\end{itemize}

\section{Examples: Future Applications}
\begin{enumerate}

\item \textbf{Same-Side (Ban Valid): A Noose on an Election Clerk's Porch.} Responding to a wave of threats against local election workers, a State enacts a statute making it a felony ``to hang a noose on the property of another or in a public place with the intent to intimidate any person,'' and adds: ``The display of a noose in public view shall be prima facie evidence of intent to intimidate.'' Defendant A hangs a noose on a Black election clerk's porch at night. Under \textit{Black} Parts I--III, the core ban on hanging a noose with intent to intimidate is constitutional: like cross burning, a noose has a uniquely virulent history as a signal of impending racial and vigilante violence, and because the ban applies regardless of the victim's race or viewpoint, it fits \textit{R.A.V.}'s virulence exception. If the jury is instructed to find intent independently from the context and is given no prima facie instruction, the core ban can support a conviction if the provision is severable or has a valid limiting construction under state law.

\item \textbf{Same-Side (Ban Valid): An Effigy Burned in a Neighbor's Driveway.} A State makes it a crime to burn any object on another person's property with intent to place that person in fear of bodily harm. After a dispute, a man burns an effigy of his neighbor in the neighbor's driveway while shouting that the neighbor will be next. The statute does not single out any message or symbol; it targets intimidation, which \textit{Black} identifies as a type of true threat directed at a person with intent to place the victim in fear of bodily harm or death. Because the State must prove that intent from the facts, rather than presuming it from the act, the prosecution fits within the unprotected category and raises none of the \textit{R.A.V.} concerns that a content-selective ban would.

\item \textbf{Opposite-Side (Prima Facie Clause Invalid): A Noose in a Protest Sculpture.} Under the same noose statute, Defendant B, a political artist, hangs a noose on a public plaza sculpture to protest capital punishment, and the prosecution relies on the prima facie clause rather than proof of intent. If the jury is instructed that the public display alone suffices to infer intent, Defendant B has the same First Amendment objection that succeeded as to Black's conviction. Four Justices treated the clause as facially invalid under that construction. The Court left open whether a different construction might save it.

\item \textbf{Opposite-Side (Statute or Prosecution Invalid): A Viewpoint-Selective Ban and an Objective Standard.} Now remove the prima facie clause and alter the statute so it bans ``burning a cross or hanging a noose \textit{to intimidate any person on the basis of race, religion, or political party},'' or alternatively prosecutes Defendant B under a general threat statute using an \textit{objective reasonable-person standard} with no subjective mens rea. Under the first variation, the constitutional result depends on what the list does. If it selects punishable threats by the viewpoint or topic they express, it raises the \textit{R.A.V.} problem. If it turns on bias in selecting the victim of independently punishable conduct, \textit{Wisconsin v.\ Mitchell} supplies a different analysis. Under the second variation, prosecuting Defendant B without proving at least subjective \textit{recklessness} that others would view the display as threatening violence violates \textit{Counterman v. Colorado} (2023).

\item \textbf{Fence-Sitter (Close Call): A Burning-Cross Image Posted Online.} During a local election, a man posts an image of a burning cross on social media and tags the account of a Black candidate, adding no words. He is prosecuted under a statute banning communications made with intent to intimidate. He says the post expressed his political views to his followers; the candidate says she understood it as a threat against her. Under \textit{Black}, a burning cross directed at a person can be a true threat, but the same symbol can also be political expression, and context decides which. Tagging the candidate points toward targeting her; the absence of any physical presence or explicit statement points the other way. Whether the State can prove intent to intimidate is genuinely uncertain.

\end{enumerate}

\section{Critique}
\textbf{Tension with \textit{R.A.V.}:} Justice Souter highlights a real doctrinal strain: if a symbol (the burning cross) is historically fused with a specific political ideology (white supremacy), does upholding a cross-burning ban under \textit{R.A.V.}'s ``virulence'' exception collapse the line between targeting harm and targeting a disfavored viewpoint?

\textbf{Procedural Overbreadth vs. Substantive Reach:} Justice Scalia's critique exposes a puzzle in First Amendment facial-challenge doctrine---whether an evidentiary inference in a jury instruction can render an otherwise valid criminal prohibition facially void even when a state supreme court has not yet construed it.

\section{Key Quotations}
\begin{quote}
````True threats'' encompass those statements where the speaker means to communicate a serious expression of an intent to commit an act of unlawful violence to a particular individual or group of individuals. \dots Intimidation in the constitutionally proscribable sense of the word is a type of true threat, where a speaker directs a threat to a person or group of persons with the intent of placing the victim in fear of bodily harm or death.''
\end{quote}
\noindent --- Justice O'Connor, 538 U.S. at 359--360 (\href{https://supreme.justia.com/cases/federal/us/538/343/}{Justia}).

\begin{quote}
``The First Amendment permits Virginia to outlaw cross burnings done with the intent to intimidate because burning a cross is a particularly virulent form of intimidation.''
\end{quote}
\noindent --- Justice O'Connor, 538 U.S. at 363 (\href{https://supreme.justia.com/cases/federal/us/538/343/}{Justia}).

\begin{quote}
``The act of burning a cross may mean that a person is engaging in constitutionally proscribable intimidation. But that same act may mean only that the person is engaged in core political speech. \dots The prima facie evidence provision in this case ignores all of the contextual factors that are necessary to decide whether a particular cross burning is intended to intimidate. The First Amendment does not permit such a shortcut.''
\end{quote}
\noindent --- Justice O'Connor (plurality), 538 U.S. at 365, 367 (\href{https://supreme.justia.com/cases/federal/us/538/343/}{Justia}).

\begin{quote}
``Accordingly, this statute prohibits only conduct, not expression. And, just as one cannot burn down someone's house to make a political point and then seek refuge in the First Amendment, those who hate cannot terrorize and intimidate to make their point.''
\end{quote}
\noindent --- Justice Thomas, dissenting, 538 U.S. at 394 (\href{https://supreme.justia.com/cases/federal/us/538/343/}{Justia}).


\end{document}
